% Options for packages loaded elsewhere
\PassOptionsToPackage{unicode}{hyperref}
\PassOptionsToPackage{hyphens}{url}
%
\documentclass[
]{article}
\usepackage{amsmath,amssymb}
\usepackage{iftex}
\ifPDFTeX
  \usepackage[T1]{fontenc}
  \usepackage[utf8]{inputenc}
  \usepackage{textcomp} % provide euro and other symbols
\else % if luatex or xetex
  \usepackage{unicode-math} % this also loads fontspec
  \defaultfontfeatures{Scale=MatchLowercase}
  \defaultfontfeatures[\rmfamily]{Ligatures=TeX,Scale=1}
\fi
\usepackage{lmodern}
\ifPDFTeX\else
  % xetex/luatex font selection
\fi
% Use upquote if available, for straight quotes in verbatim environments
\IfFileExists{upquote.sty}{\usepackage{upquote}}{}
\IfFileExists{microtype.sty}{% use microtype if available
  \usepackage[]{microtype}
  \UseMicrotypeSet[protrusion]{basicmath} % disable protrusion for tt fonts
}{}
\makeatletter
\@ifundefined{KOMAClassName}{% if non-KOMA class
  \IfFileExists{parskip.sty}{%
    \usepackage{parskip}
  }{% else
    \setlength{\parindent}{0pt}
    \setlength{\parskip}{6pt plus 2pt minus 1pt}}
}{% if KOMA class
  \KOMAoptions{parskip=half}}
\makeatother
\usepackage{xcolor}
\setlength{\emergencystretch}{3em} % prevent overfull lines
\providecommand{\tightlist}{%
  \setlength{\itemsep}{0pt}\setlength{\parskip}{0pt}}
\setcounter{secnumdepth}{-\maxdimen} % remove section numbering
\ifLuaTeX
  \usepackage{selnolig}  % disable illegal ligatures
\fi
\IfFileExists{bookmark.sty}{\usepackage{bookmark}}{\usepackage{hyperref}}
\IfFileExists{xurl.sty}{\usepackage{xurl}}{} % add URL line breaks if available
\urlstyle{same}
\hypersetup{
  hidelinks,
  pdfcreator={LaTeX via pandoc}}

\author{}
\date{}

\begin{document}

\newcommand{\ba}{\begin{aligned}}
\newcommand{\ea}{\end{aligned}}
\newcommand{\bm}{\begin{bmatrix}}
\newcommand{\em}{\end{bmatrix}}
\newcommand{\bc}{\begin{cases}}
\newcommand{\ec}{\end{cases}}

\newcommand{\empty}{\varnothing}
\newcommand{\se}{\subseteq}
\newcommand{\ve}{\varepsilon}
\newcommand{\ep}{\epsilon}
\newcommand{\s}{\sigma}

\newcommand{\0}{\textbf 0}
\newcommand{\1}{\textbf 1}
\newcommand{\I}{\textbf I}

\newcommand{\X}{\textbf X}
\newcommand{\x}{\textbf x}
\newcommand{\y}{\textbf y}
\newcommand{\b}{\textbf b}


\newcommand{\P}{\mathcal P}
\newcommand{\S}{\mathcal S}
\newcommand{\T}{\mathcal T}
\newcommand{\A}{\mathcal A}
\newcommand{\B}{\mathcal B}
\newcommand{\L}{\mathcal L}

\newcommand{\NN}{\mathbb N}
\newcommand{\ZZ}{\mathbb Z}
\newcommand{\RR}{\mathbb R}
\newcommand{\QQ}{\mathbb Q}
\newcommand{\CC}{\mathbb C}
\newcommand{\FF}{\mathbb F}

\newcommand{\c}{\backslash}
\newcommand{\sq}{\sqrt}

\newcommand{\too}{\Longrightarrow}
\newcommand{\dto}{\overset{\text{d}}{\to}}
\newcommand{\ae}{\overset{\text{a.e.}}{\to}}
\newcommand{\asim}{\overset{\text a}{\sim}}
\newcommand{\st}{\text{ s.t. }}
\newcommand{\plim}{\text{plim }}
\newcommand{\if}{\text{ if }}

\newcommand{\uinf}[1][k]{\bigcup_{{#1} = 1}^{\infty}}
\newcommand{\sinf}[1][k]{\sum_{{#1} = 1}^{\infty}}
\newcommand{\iinf}[1][k]{\bigcap_{{#1} = 1}^{\infty}}
\newcommand{\seq}[1][x]{{#1}_1,{#1}_2,\cdots,{#1}_{n}}
\newcommand{\seqadd}[1][x]{{#1}_1+{#1}_2+\cdots+{#1}_{n}}

\newcommand{\var}[2][]{\text{Var}_{#1}\br{#2}}
\newcommand{\e}[2][]{\text{E}_{#1}\br{#2}}
\newcommand{\cov}[3][]{\text{Cov}_{#1}\br{#2,#3}}
\newcommand{\estvar}[1]{\text{Est.Var}\br{#1}}
\newcommand{\estcov}[2]{\text{Est.Cov}\br{#1,#2}}
\newcommand{\asyvar}[1]{\text{Asy.Var}\br{#1}}
\newcommand{\asycov}[2]{\text{Asy.Cov}\br{#1,#2}}
\newcommand{\estasyvar}[1]{\text{Est.Asy.Var}\br{#1}}

\newcommand{\br}[1]{\left[#1\right]}
\newcommand{\pr}[1]{\left(#1\right)}
\newcommand{\brace}[1]{\left\{#1\right\}}
\newcommand{\dev}[2][i]{({#2}_{#1}-\bar{#2})}
\newcommand{\inv}[1]{{#1}^{-1}}
\newcommand{\add}[1][i]{\sum_{#1 = 1}^n}
\newcommand{\ply}[1][i]{\prod_{#1 = 1}^n}
\newcommand{\om}[1]{\mu^*\pr{#1}}
\newcommand{\lm}[1]{\mu\pr{#1}}
\newcommand{\d}{\text{ d}}
\newcommand{\dm}{\text{ d}\mu}
\newcommand{\dl}{\text{ d}\lambda}
\newcommand{\norm}[2][1]{\|{#2}\|_{#1}}
\newcommand{\limit}[1][k]{\lim_{{#1}\to\infty}}

\hypertarget{integration}{%
\section{Integration}\label{integration}}

\hypertarget{integration-with-respect-to-a-measure}{%
\subsection{Integration with Respect to a
Measure}\label{integration-with-respect-to-a-measure}}

\hypertarget{integration-of-nonnegative-functions}{%
\subsubsection{Integration of Nonnegative
Functions}\label{integration-of-nonnegative-functions}}

Through out this chapter we base our analysis on \(\pr{X,\S,\mu}\),
which is a measure space.

\textbf{Definition (\(\S\)-Partition)} An \(\S\)-partition of \(X\) is a
finite collection \(A_1,\cdots,A_m\) of disjoint sets in \(\S\) such
that

\[A_1\cup\cdots \cup A_m = X\]

\textbf{Definition (Lower Lebesgue Sum)} Suppose
\(f:X\to \br{0,\infty}\) is an \(\S\)-measurable function, and \(P\) is
an \(\S\)-partition \(\brace{A_1,\cdots, A_m}\) of \(X\). The lower
Lebesgue sum \(L\pr{f,P}\) is defined by

\[L\pr{f,P}:=\sum_{j = 1}^m\lm{A_j}\inf_{A_j}f\]

\textbf{Remarks:}

\begin{itemize}
\item
  Notice that here we assume \(f\) to be nonnegative, so
  \(\inf_{A_j} f\) always exists.
\item
  One convention is that, if any of the two terms is \(0\), the result
  is \(0\) and then neglected.
\end{itemize}

\textbf{Definition (Integral of a Nonnegative Function)} The integral of
\(f\) with respect to \(\mu\) denoted by \(\int f \d u\), is defined by

\[\int f\d u = \sup_{P}\brace{L\pr{f,P}:P \text{ is an }\S\text{-partition of }X}\]

We should be able to prove the following:

\begin{itemize}
\item
  \(\int \chi_E\d \mu = \lm{E}\).
\item
  Additivity: \(\int \pr{f+g}\d \mu = \int f\d \mu+\int g\d \mu\).
\item
  Multiplication with a constant:
  \(\int cf\d \mu = c\cdot \int f\d \mu\).
\end{itemize}

\emph{\textbf{Proposition}} If \(E\in \S\), then

\[\int\chi_E\d \mu = \lm{E}\]

\emph{\textbf{Proof}} We separate the task to prove both
\(\int\chi_E\d \mu \le \lm{E}\) and \(\int\chi_E\d \mu \ge \lm{E}\).

Consider the partition \(E,E^c\). Obviously \(E\cup E^c = X\). From
definition,

\[\ba
\int\chi_E\dm & \ge L\pr{\chi_E,\brace{E,E^c}} \\
& = \lm{E}\cdot \inf_{E}\chi_E+\lm{E^c}\cdot \inf_{E^c}\chi_{E} \\
& = \lm{E}+0 \\
& = \lm{E}
\ea\]

Pick any partition, then prove \(L\pr{\chi_E,P}\le \lm{E}\).

Consider any \(\S\)-partition \(P\) such that
\(A_1\cup\cdots \cup A_m = X\).

\[\ba
\L\pr{\chi_E,P} & = \sum_{j = 1}^m\lm{A_j}\inf_{A_j}\chi_E\\
& = \sum_{j:A_j\se E}\lm{A_j}\cdot 1\\
& = \lm{\bigcup_{j:A_j\se E}A_j}\\
& \le \lm{E}
\ea\]

Taking supremium to both sides, we end with
\(\int \chi_E\dm \le \lm{E}\), which when combined with the previous
result completes the proof.

\emph{\textbf{Proposition}} Suppose \(E_1,E_2,\cdots ,E_n\) are disjoint
sets in \(\S\). \(c_1,c_2,\cdots,c_n\in \br{0,\infty}\). Then

\[\int \pr{\sum_{k = 1}^nc_k\chi_{E_k}}\dm = \sum_{k = 1}^nc_k\lm{E_k}\]

\emph{\textbf{Proof}} When one of the \(c_1,\cdots,c_n\) is \(\infty\),
the equality is trivially true.

Let \(E_{n+1} = X\c \pr{E_1\cup\cdots \cup E_n}\) and \(c_{n+1} = 0\).
Then \(P = \brace{E_1,E_2,\cdots ,E_n,E_{n+1}}\) is an \(\S\)-partition
of \(X\). We have

\[\ba
\int \pr{\sum_{k = 1}^nc_k\chi_{E_k}}\dm & = \int \pr{\sum_{k = 1}^{n+1}c_k\chi_{E_k}}\dm \\
& \ge L\pr{\sum_{k = 1}^{n+1}c_k\chi_{E_k},\brace{E_1,\cdots,E_{n+1}}}\\
& = \sum_{k = 1}^{n+1}c_k\lm{E_k}\\
& = \sum_{k = 1}^{n}c_k\lm{E_k}
\ea\]

Suppose \(P\) is an \(\S\)-partition \(\brace{A_1,\cdots ,A_m}\) of
\(X\). Then

\[\ba
L\pr{\sum_{k = 1}^nc_k\chi_{E_k},P} & = \sum_{j = 1}^m\lm{A_j}\inf_{A_j}\sum_{k = 1}^nc_k\chi_{E_k} \\
& = \sum_{j = 1}^m\lm{A_j}\min_{\brace{i:A_j\cap E_i\ne \empty}}c_i\\
& = \sum_{j = 1}^m\pr{\sum_{k = 1}^n \lm{A_j\cap E_k}\min_{\brace{i:A_j\cap E_i\ne \empty}}c_i}\\
& \le \sum_{j = 1}^m\pr{\sum_{k = 1}^n \lm{A_j\cap E_k}c_k}\\
& = \sum_{k = 1}^nc_k\sum_{j = 1}^m\lm{A_i\cap E_k}\\
& = \sum_{k = 1}^nc_k \lm{E_k}
\ea\]

Note that the equality in the last line stands because of

\[\ba
A_1\cup\cdots \cup A_m & = X\\
\sum_{j = 1}^m\lm{A_j\cap E_k} & = \lm{\pr{A_1\cup\cdots \cup A_m}\cap E_k}\\
& = \lm{X\cap E_k}\\
& = \lm{E_k}
\ea\]

\emph{\textbf{Proposition}} If \(f\pr x\le g\pr x\) for all \(x\in X\),
then

\[\int f\dm \le \int g\dm\]

\hypertarget{monotone-convergence-theorem}{%
\subsubsection{Monotone Convergence
Theorem}\label{monotone-convergence-theorem}}

\textbf{Definition (Integrals via Finite Simple Functions)}

\[\ba
\int f\dm := \sup_{\brace{A_j,c_j}} \bigg\{\sum_{j = 1}^m c_j\lm{A_j}: & A_1,\cdots ,A_m\text{ are disjoint sets of }X,\\ 
& c_1,\cdots ,c_m \in \br{0,\infty},\\
& f\pr{x}\ge \sum_{j = 1}^m c_j\chi_{A_j}\pr{x},\forall x\in X\bigg\}
\ea\]

\emph{\textbf{Monotone Convergence Theorem}} Let
\(0\le f_1\le f_2\le \cdots\) is an increasing sequence of
\(\S\)-measurable functions. Define \(f: X\to \br{0,\infty}\) as
\(f\pr{x}:=\lim_{k\to \infty}f_k\pr{x}\). Then

\[\lim_{k\to\infty} \int f_k\dm = \int f\dm\]

\emph{\textbf{Proof}} Since \(f_1\le f_2\le \cdots \le f\),

\[\ba
& \int f_1\dm \le \int f_2\dm \le \cdots \le \int f\dm\\
\too & \lim_{k\to \infty}\int f_k\dm \le \int f\dm
\ea\]

Now we are to focus on the other direction. Using the equivalent
definition of integration, suppose \(A_1,\cdots ,A_m\) are disjoint sets
of \(X\), and \(c_1,\cdots ,c_m\in \br{0,\infty}\) such that

\[f\pr{x}\ge \sum_{j = 1}^mc_j\chi_{A_j}\pr{x},\forall x\in X\]

Let \(t\in \pr{0,1}\). For \(k\in \Z_+\), let

\[E_k = \brace{x\in X: f_k\pr{x}\ge t \sum_{j = 1}^m c_j\chi_{A_j}\pr x}\]

Then \(E_1\se E_2\se \cdots\) is an increasing sequence whose union
\(\uinf E_k = X\). Thus

\[\lim_{k\to\infty}\lm{A_j\cap E_k} = \lm{A_j\cap \pr{\uinf E_k}} = \lm{A_j},\forall j\in \brace{1,2,\cdots,m}\]

Moreover, from the construction of \(E_k\),

\[f_k\pr{x}\ge t\sum_{j = 1}^m c_j\chi_{A_j\cap E_k}\pr{x},\forall x\in X\]

Hence,

\[\int f_k\dm \ge t\sum_{j = 1}^mc_j\lm{A_j\cap E_k}\]

Let \(k\to \infty\), we have

\[\lim_{k \to \infty}\int f_k\dm \ge t\sum_{j = 1}^mc_j\lm{A_j}\]

Take \(t\to 1\), we have

\[\lim_{k \to \infty}\int f_k\dm \ge \sum_{j = 1}^mc_j\lm{A_j}\]

\textbf{Remarks:} Why only take \(t\in \pr{0,1}\) instead of
\(t\in (0,1]\). If \(t = 1\), we may not have \(\uinf E_k = X\).

\begin{center}\rule{0.5\linewidth}{0.5pt}\end{center}

Next we will deploy Monotone Convergence Theorem to prove that the
integration has the property of additivity, as is desired.

A simple function may have different representations, but the
corresponding integrals are of the same value, which is formalized in
the following proposition.

\emph{\textbf{Proposition (Integral-Type Sums for Simple Functions)}}
Suppose \(\pr{X,\S,\mu}\) is a measure space. Suppose
\(a_1,a_2,\cdots,a_m,b_1,b_2,\cdots ,b_n\in [0,\infty)\), and
\(A_1,A_2,\cdots,A_m,B_1,B_2,\cdots ,B_n\se\S\). Further suppose

\[\sum_{i = 1}^ma_i\chi_{A_i}\pr{x} = \sum_{j = 1}^nb_j\chi_{B_j}\pr{x},\forall x\in X\]

Then

\[\sum_{i = 1}^ma_i\lm{A_i} = \sum_{j = 1}^nb_j\lm{B_j}\]

\textbf{Remarks}:

\begin{itemize}
\item
  Here it is not required that \(\brace{a_i},\brace{b_i}\) are distinct
  or \(\brace{A_k},\brace{B_k}\) are disjoint. From here we do not need
  to care about multiplicity about representations of simple functions.
\item
  The standard representation for a simple function:
  \(\sum_{i = 1}^m a_i\chi_{A_i}\), where \(a_1,a_2,\cdots ,a_m\) are
  distinct and \(A_1,A_2,\cdots ,A_m\) are disjoint.
\end{itemize}

\emph{\textbf{Proposition (Additivity of Integration)}} If \(f,g\ge 0\)
are \(\S\)-measurable, then

\[\int\pr{f+g}\dm = \int f\dm+\int g\dm\]

\emph{\textbf{Proof}} Let \(\brace{f_k}\) converges pointwise to \(f\)
in an increasing order, so does \(\brace{g_k}\) to \(g\). By MCT
(applied to \(\brace{f_k+g_k}\)),

\[\ba
\int\pr{f+g}\dm & = \lim_{k\to\infty}\int \pr{f_k+g_k}\dm\\
& = \lim_{k\to \infty}\pr{\int f_k\dm +\int g_k\dm}\\
& = \int f\dm +\int g\dm
\ea\]

where the equality in the second line stands because both \(f_k\) and
\(g_k\) are simple functions, and the last one holds when applying MCT
to \(\brace{f_k}\) and \(\brace{g_k}\) respectively.

\hypertarget{integration-of-real-valued-functions}{%
\subsubsection{Integration of Real-Valued
Functions}\label{integration-of-real-valued-functions}}

\textbf{Definition (\(f^+\), \(f^-\))} Suppose
\(f:X\to \br{-\infty,\infty}\) is \(\S\)-measurable. Define
\(f^+:X\to \br{0,\infty}\) and \(f^-:X\to \br{0,\infty}\) by

\[\ba
f^+\pr{x}& = \bc
f\pr{x} & \if f\pr{x}\ge 0\\
0 & \if f\pr{x}<0
\ec\\
f^-\pr{x}& = \bc
0& \if f\pr{x}\ge 0\\
-f\pr{x} & \if f\pr{x}<0
\ec
\ea\]

By definition, we see that

\[\ba
f\pr{x} & = f^+\pr{x}-f^-\pr{x}\\
\abs{f}\pr{x} & = \abs{f^+}\pr{x}+\abs{f^-}\pr{x}
\ea\]

\textbf{Definition (Integral of a Real-Valued Function)} Suppose
\(f:X\to \br{-\infty,\infty}\) is an \(\S\)-measurable function such
that at least one \(\int f^+\dm\) and \(\int f^-\dm\) is finite. The
integral of \(f\) is

\[\int f\dm := \int f^+\dm -\int f^-\dm\]

\emph{\textbf{Proposition (Properties of Integration)}} Suppose
\(f:X\to \br{-\infty,\infty}\), and \(g:X\to \br{-\infty,\infty}\).

\begin{itemize}
\item
  \(\int cf\dm = c\int f\dm\), if \(c\in \R\).
\item
  \(\int\pr{f+g}\dm = \int f\dm + \int g\dm\), if
  \(\int \abs{f}\dm<\infty\) and \(\int \abs{g}\dm<\infty\).
\item
  \(\int f\dm \le \int g\dm\), if \(f\pr{x}\le g\pr x\) for all
  \(x\in X\).
\item
  Triangle inequality: \(\abs{\int f\dm}\le \int \abs{f}\dm\).
\end{itemize}

\hypertarget{limits-of-integrals--integrals-of-limits}{%
\subsection{Limits of Integrals \& Integrals of
Limits}\label{limits-of-integrals--integrals-of-limits}}

\hypertarget{bounded-convergence-theorem}{%
\subsubsection{Bounded Convergence
Theorem}\label{bounded-convergence-theorem}}

\textbf{Definition (Integration on a Subset)} Suppose \(\pr{X,\S,\mu}\)
is a measure space and \(E\in \S\). Let \(f:X\to \br{-\infty,\infty}\)
be a \(\S\)-measurable function. The integral of \(f\) over \(E\) is
defined by

\[\int_Ef\dm:=\int f\cdot\chi_E\dm\]

The special case is that
\(\int_Xf\dm = \int f\cdot \chi_X\dm = \int f\dm\). So when we refer to
integration over the whole set \(X\), we can simply write as
\(\int f\dm\) and omit the denotation of \(X\).

\emph{\textbf{Proposition (Bounding an Integral)}} Suppose
\(\pr{X,\S,\mu}\) is a measure space, \(E\in \S\), and
\(f:X\to\br{-\infty,\infty}\) is a function such that \(\int_Ef\dm\) is
defined. Then

\[\abs{\int_Ef\dm}\le \lm{E}\cdot\sup_E\abs{f}\]

\emph{\textbf{Proof}}

\[\ba
\abs{\int_Ef\dm} & = \abs{\int f\cdot\chi_E\dm}\\
& \le \int \abs{f}\cdot\chi_E\dm \\
& \le \int \sup_E\abs{f}\cdot\chi_E\dm\\
& = \lm{E}\cdot\sup_E\abs{f}
\ea\]

\emph{\textbf{Bounded Convergence Theorem}} Suppose \(\pr{X,\S,\mu}\) is
a measure space with \(\lm{X}<\infty\). Suppose a sequence of
\(\S\)-measurable functions \(f_1,f_2,\cdots\) converges pointwise to
\(f\). Suppose there exists \(c\in\pr{0,\infty}\) such that
\(\abs{f_k\pr{x}}\le c\), for all \(x\in X,k\ge 1\). Then

\[\lim_{k\to\infty}\int f_k\dm = \int f\dm\]

\emph{\textbf{Proof}} Fix \(\ve>0\). By Egorov\textquotesingle s
theorem, there exists \(E\in \S\) such that

\[\lm{X\c E}\le \frac{\ve}{4c}\]

and \(f_1,f_2,\cdots\) converges to \(f\) uniformly on \(E\). Then

\[\ba
\abs{\int f_k\dm-\int f\dm} & = \abs{\int_{X\c E}f_k\dm-\int_{X\c E}f\dm+\int_Ef_k\dm-\int_Ef\dm}\\
& \le \abs{\int_{X\c E}\pr{f_k-f}\dm}+\abs{\int_E\pr{f_k-f}\dm}\\
& \le \lm{X\c E}\sup_{X\c E}\abs{f_k-f}+\lm{E}\sup_{E}\abs{f_k-f}\\
& \le \frac{\ve}{4c}\cdot 2c+\lm{X}\cdot \frac{\ve}{2\pr{\lm{X}+1}}\\
& \le \frac{\ve}{2}+\frac{\ve}2\\
& = \ve
\ea\]

Hence, \(\abs{\int f_k\dm-\int f\dm}\le \ve\) when \(k\) is big enough.
This implies that

\[\lim_{k\to \infty}\int f_k\dm = \int f\dm\]

\textbf{Remarks}: Here it is in the last approximation of
\(\sup_{E}\abs{f_k-f}\) that we used the definition (or, property) of
uniform convergence. Egorov\textquotesingle s theorem is a powerful
tools in proofs that involves interchanging limits and integrals.

\hypertarget{dominated-convergence-theorem}{%
\subsubsection{Dominated Convergence
Theorem}\label{dominated-convergence-theorem}}

Suppose \(\pr{X,\S,\mu}\) is a measure space. If
\(f,g:X\to\br{-\infty,\infty}\) are \(\S\)-measurable functions and
\(\lm{\brace{x\in X:f\pr{x}\ne g\pr{x}}} = 0\), then the definition of
an integral implies that \(\int f\dm = \int g\dm\) (or both are
undefined). This is true because what happens on a set of measure \(0\)
does not matter generally.

\textbf{Definition (\(\mu\)-Almost Every)} Suppose \(\pr{X,\S,\mu}\) is
a measure space. A set \(E\in \S\) is said to contain \(\mu\)-almost
every element of \(X\) if

\[\lm{X\c E} = 0\]

\textbf{Example} Almost every real number is irrational (with respect to
the usual Lebesgue measure on \(\R\)) because \(\om{\Q} = 0\).

\textbf{Remarks}

\begin{itemize}
\item
  For two functions \(f\) and \(g\), if
  \(\lm{\brace{x\in X:f\pr{x}\ne g\pr{x}}} = 0\), we usually say
  \(f = g\) (\(\mu\)-)almost everywhere (a.e.).
\item
  Theorems about integrals can almost always be relaxed so that the
  hypothesis apply only almost everywhere instead of everywhere. One
  example is \(\lim_{k\to\infty}f_k\pr{x} = f\pr{x}\) for all \(x\in X\)
  in Bounded Convergence Theorem.
\end{itemize}

\emph{\textbf{Proposition (Integrals on Small Sets are Small)}} Suppose
\(\pr{X,\S,\mu}\) is a measure space. Suppose \(g:X\to\br{0,\infty}\) is
\(\S\)-measurable and \(\int g\dm <\infty\). Then for every \(\ve>0\),
there exists \(\delta>0\) such that if \(\lm{B}<\delta\), then
\(\int_{B}g\dm<\ve\).

\emph{\textbf{Proof}} Fix \(\ve>0\). Construct \(h:X\to [0,\infty)\) to
be a simple \(\S\)-measurable function such that \(0\le h\le g\) and
\(\int g\dm-\int f\dm \le \frac{\ve}2\). (The existence can be show via
the equivalent definition of Lebesgue integral.)

Let \(H = \max\brace{h\pr{x}:x\in X}\) and let \(\delta>0\) be such that
\(H\cdot \delta \le \frac{\ve}2\).

\[\ba
\int_Bg\dm & = \pr{\int_B g\dm -\int_Bh\dm}+\int_Bh\dm\\
& \le \frac{\ve}2+\lm{B}\cdot\sup_B h\\
& \le \frac{\ve}2+\delta\cdot H \\
& \le \frac{\ve}2+\frac{\ve}2\\
& = \ve
\ea\]

\emph{\textbf{Proposition (Integrable Functions Live Mostly on Sets of
Finite Measure)}} Let \(g:X\to \br{0,\infty}\) and \(\int g\dm<\infty\).
Then for every \(\ve>0\), there exists \(E\in \S\) such that
\(\lm{E}<\infty\) and \(\int_{X\c E}g\dm\le \ve\).

\emph{\textbf{Proof}} Let \(\ve>0\). Let \(P\) be an \(\S\)-partition
\(A_1,\cdots ,A_m\) of \(X\) such that

\[\int g\dm \le\ve +L\pr{g,P}\]

Let \(E\) be the union of those \(A_i\pr{i = 1,\cdots,m}\) such that
\(\inf_{A_i}g>0\). We can safely claim that \(\lm{E}<\infty\); otherwise
we would have \(L\pr{g,P} = \infty\), which contradicts the hypothesis
that \(\int g\dm<\infty\). Now

\[\ba
\int_{X\c E}g\dm & = \int g\dm -\int\chi_{E}\cdot g\dm\\
& \le \ve+L\pr{g,P}-L\pr{\chi_{E}\cdot g,P} \\
& = \ve
\ea\]

Note that the second line follows from our construction and the
definition of the integral, and the last line holds because
\(\inf_{A_j}g = 0\) for each \(A_j\) not contained in \(E\).

\emph{\textbf{Dominated Convergence Theorem}} Suppose \(\pr{X,\S,\mu}\)
is a measure space. Let \(f:X\to \br{-\infty,\infty}\) and
\(f_1,f_2,\cdots:X\to \br{-\infty,\infty}\) be \(\S\)-measurable
functions such that \(\lim_{k\to\infty}f_k\pr{x} = f\pr{x}\) for almost
every \(x\in X\). If there exists an \(\S\)-measurable function
\(g:X\to\br{0,\infty}\) such that \(\int g\dm<\infty\) and
\(\abs{f_k\pr{x}}\le g\pr{x}\) for all \(k>0\) and almost every
\(x\in X\), then

\[\lim_{k\to\infty} \int f_k\dm = \int f\dm\]

\emph{\textbf{Proof}} Suppose \(g:X\to\br{0,\infty}\) satisfies the
hypotheses of this theorem. For any set \(E\in\S\), we have the
following estimates

\[\ba
\abs{\int f_k\dm-\int f\dm} & = \abs{\int_{X\c E}f_k\dm-\int_{X\c E}f\dm+\int_Ef_k\dm-\int_Ef\dm}\\
& \le \int_{X\c E}\abs{f_k}\dm +\int_{X\c E}\abs{f}\dm+\abs{\int_E\pr{f_k-f}\dm}\\
& \le 2\int_{X\c E}g\dm+\abs{\int_E\pr{f_k-k}\dm}
\ea\]

Assume first for now that \(\lm{X}<\infty\). Let \(\ve>0\). From the
proposition that integrals on small sets are small, there exists
\(\delta>0\) such that \(\int_Bg\dm\le \frac{\ve}4\) for every set
\(B\in \S\) such that \(\lm{B}<\ve\). By Egorov\textquotesingle s
theorem (since \(\lm{X}<\infty\)), there exists \(E\in \S\) such that

\[\lm{X\c E}\le \delta\]

and \(f_1,f_2,\cdots\) converges uniformly on \(E\). We then have

\[\ba
\abs{\int f_k\dm-\int f\dm } & \le 2\int_{X\c E}g\dm+\abs{\int_E\pr{f_k-f}\dm}\\
& \le 2\cdot \frac{\ve}4+\lm{E}\cdot \sup_E\abs{f_k-f}\\
& \le \frac{\ve}2+\lm{E}\cdot \frac{2\ve}{\lm{E}}\\
& \le \frac{\ve}2+\frac{\ve}2\\
& = \ve
\ea\]

Then consider the case where \(\lm{X} = \infty\). Let \(\ve>0\). By
proposition 3, there exists \(E\in \S\) such that \(\lm{E}<\infty\) and
\(\int_{X\c E}\dm \le \frac{\ve}4\). Then we have

\[\ba
\abs{\int f_k\dm-\int f\dm} & \le 2\int_{X\c E}g\dm+\abs{\int_E\pr{f_k-f}\dm}\\
& \le 2\cdot \frac{\ve}4+\abs{\int_E\pr{f_k-k}\dm}
\ea\]

By the first case as applied to the sequence \(f_1|_E,f_2|_E,\cdots\).
the last term on the right is less than \(\frac{\ve}2\) for all
sufficiently large \(k\). Thus
\(\lim_{k\to\infty}\int f_k\dm = \int f\dm\), completing the proof of
the second case.

\textbf{Remarks:}

\begin{itemize}
\item
  Advantages of DCT

  \begin{itemize}
  \item
    We do not require the functions to be nonnegative.
  \item
    We do not require the sequence to be increasing.
  \item
    We do not require the measure space to be finite.
  \item
    We do not require the sequence of functions to be uniformly bounded.
  \end{itemize}

  All these hypotheses are replaced only by a requirement that the
  sequence of functions is pointwise bounded by a function with a finite
  integral.
\item
  Bounded Convergence Theorem follows immediately from the result of
  Dominated Convergence Theorem. Simply take \(g\) to be an appropriate
  constant function and use the hypothesis in the Bounded Convergence
  Theorem that \(\lm{X}<\infty\).
\end{itemize}

\hypertarget{relationship-between-lebesgue-integral-and-riemann-integral}{%
\subsubsection{Relationship Between Lebesgue Integral and Riemann
Integral}\label{relationship-between-lebesgue-integral-and-riemann-integral}}

\emph{\textbf{Theorem (Riemann Integrable is Equivalent to Continuous
Almost Everywhere)}} Suppose \(a<b\) and \(f:\br{a,b}\to \R\) is a
bounded function. Then \(f\) is Riemann integrable if and only if
\(\om{\brace{x\in\br{a,b}:f\text{ is not continuous at }x}} = 0\).
Moreover, if \(f\) is Riemann integrable,

\[\int_a^bf\d x = \int_{\br{a,b}}f\d \lambda\]

where \(\lambda\) is the Lebesgue measure.

\end{document}
